% =======================================================
% АССОЦИАТИВНЫЕ КОНТЕЙНЕРЫ
% =======================================================

\section{STL: ассоциативные контейнеры}

% -------------------------------------------------------
\subsection{Обзор}

\textbf{Ассоциативные контейнеры} хранят элементы, доступные по \textbf{ключу}, а
не по позиции. Делятся на две группы:

\begin{center}
\begin{tabular}{lll}
\toprule
\textbf{Группа} & \textbf{Контейнеры} & \textbf{Сложность поиска} \\
\midrule
Упорядоченные (дерево) & \texttt{map}, \texttt{set}, \texttt{multimap}, \texttt{multiset} & $O(\log n)$ \\
\addlinespace
Неупорядоченные (хеш) & \texttt{unordered\_map}, \texttt{unordered\_set}, \dots & $O(1)$ в среднем \\
\bottomrule
\end{tabular}
\end{center}

\subsection{\texttt{std::set} --- упорядоченное множество}

Хранит \textbf{уникальные} ключи в отсортированном порядке. Реализован
красно-чёрным деревом.

\begin{lstlisting}
#include <set>
std::set<int> s = {3, 1, 4, 1, 5};   // дубликат 1 игнорируется
// хранится отсортированно: 1 3 4 5
s.insert(2);
s.count(3);        // 1 (есть) или 0 (нет)
s.contains(3);     // C++20: bool
s.erase(4);
for (int x : s) std::cout << x;   // обход в порядке возрастания
\end{lstlisting}

\subsection{\texttt{std::map} --- упорядоченный словарь}

Хранит пары \textbf{ключ--значение} с уникальными ключами, отсортированными по
ключу.

\begin{lstlisting}
#include <map>
std::map<std::string, int> ages;
ages["Alice"] = 30;          // вставка/изменение через []
ages["Bob"]   = 25;
ages.insert({"Carol", 28});

for (const auto& [name, age] : ages)   // structured bindings (C++17)
    std::cout << name << ": " << age << '\n';   // в порядке ключей
\end{lstlisting}

\subsubsection{Осторожно с \texttt{operator[]}}

\texttt{operator[]} \textbf{создаёт} элемент со значением по умолчанию, если ключа
нет --- это меняет map и не работает на \texttt{const map}:

\begin{lstlisting}
std::map<std::string, int> m;
int x = m["missing"];   // СОЗДАЁТ запись {"missing", 0}!

// безопасные способы чтения:
auto it = m.find("key");
if (it != m.end()) use(it->second);
int y = m.at("key");    // бросает std::out_of_range, если ключа нет
\end{lstlisting}

\subsubsection{Эффективная вставка}

\begin{lstlisting}
m.insert_or_assign("key", 5);       // C++17: вставить или перезаписать
m.try_emplace("key", expensive());  // C++17: не вычисляет value, если ключ есть
auto [it, inserted] = m.emplace("k", 1);  // inserted == был ли вставлен
\end{lstlisting}

\subsection{\texttt{multiset} и \texttt{multimap}}

Допускают \textbf{дубликаты} ключей. У \texttt{multimap} нет \texttt{operator[]}.
Все элементы с данным ключом получают через \texttt{equal\_range}:

\begin{lstlisting}
std::multimap<std::string, int> mm;
mm.insert({"a", 1});
mm.insert({"a", 2});       // дубликат ключа разрешён
auto [begin, end] = mm.equal_range("a");
for (auto it = begin; it != end; ++it) std::cout << it->second;
\end{lstlisting}

\subsection{Неупорядоченные контейнеры (хеш-таблицы)}

\texttt{unordered\_set} / \texttt{unordered\_map} хранят элементы в хеш-таблице:
поиск/вставка в среднем $O(1)$, но порядок не определён.

\begin{lstlisting}
#include <unordered_map>
std::unordered_map<std::string, int> h;
h["one"] = 1;       // O(1) в среднем
h.count("one");
\end{lstlisting}

\subsubsection{Устройство хеш-таблицы}

Элементы раскладываются по \textbf{корзинам (buckets)} по значению хеш-функции.
При совпадении хешей (\textbf{коллизии}) элементы в одной корзине образуют цепочку.
\textbf{Load factor} = число элементов / число корзин; при превышении порога
происходит \textbf{rehash} --- перестроение таблицы.

\begin{lstlisting}
h.load_factor();          // текущая загрузка
h.max_load_factor(0.7);   // порог рехеша
h.reserve(1000);          // заранее выделить корзины, избежать рехешей
\end{lstlisting}

\subsubsection{Свой хеш для пользовательского типа}

\begin{lstlisting}
struct Point { int x, y; };
struct PointHash {
    size_t operator()(const Point& p) const {
        return std::hash<int>{}(p.x) ^ (std::hash<int>{}(p.y) << 1);
    }
};
std::unordered_map<Point, int, PointHash> grid;
\end{lstlisting}

\subsection{Упорядоченные vs неупорядоченные: что выбрать}

\begin{center}
\begin{tabular}{lp{4.3cm}p{4.3cm}}
\toprule
& \textbf{map/set (дерево)} & \textbf{unordered (хеш)} \\
\midrule
Поиск/вставка & $O(\log n)$ & $O(1)$ в среднем, $O(n)$ худший \\
\addlinespace
Порядок обхода & отсортирован & произвольный \\
\addlinespace
Требование к ключу & \texttt{operator<} & хеш + \texttt{operator==} \\
\addlinespace
Range-запросы & да (\texttt{lower\_bound}) & нет \\
\addlinespace
Инвалидация итераторов при вставке & нет & да (при рехеше) \\
\bottomrule
\end{tabular}
\end{center}

\textbf{Правило:} нужен порядок или диапазонные запросы --- \texttt{map}/\texttt{set};
нужна максимальная скорость поиска по точному ключу --- \texttt{unordered\_*}.

\subsection{Кастомный компаратор для упорядоченных контейнеров}

\begin{lstlisting}
std::set<int, std::greater<int>> desc;   // по убыванию
// свой компаратор:
struct ByLength {
    bool operator()(const std::string& a, const std::string& b) const {
        return a.size() < b.size();
    }
};
std::set<std::string, ByLength> by_len;
\end{lstlisting}
